(1) 考虑稳态、二维、不可压、恒密度 \(\rho>0\) 的反向磁场 \(B_x\simeq\pm B_0\) 在长薄电流片中重联。电流片半长 \(L\)，半厚 \(\delta\ll L\)，等效磁扩散率 \(\eta_m>0\)（单位 \(\mathrm{m^2/s}\)）。忽略压力的主导梯度以外的细节，面外Ohm定律写为 \(E_z=-({\boldsymbol v}\times{\boldsymbol B})_z+\mu_0\eta_mJ_z\)。推导外部入流与内部电阻层电场匹配及质量守恒的主导阶关系。
(2) 用上游磁能密度与出流动压量级匹配定义 \(v_A=B_0/\sqrt{\mu_0\rho}\)，与上一问联立推导 \(\delta, v_{\rm in}\) 与Lundquist数 \(S=Lv_A/\eta_m\) 的缩放，写出磁通重联电场大小。
(3) 把入流动压相对磁压、片层纵横比以及重联率分别写为 \(S\) 的幂律，说明高 \(S\) 时为何经典重联变慢；验证所有公式量纲及 \(\eta_m\to0^+\) 的极限。
(4) 对单位面外深度，用两侧入流Poynting通量估计磁能流入率。利用喷流横截面及 \(v_{\rm out}\sim v_A\) 估算动能通量，比较其数量级。解释为何仅靠薄片内常数电流密度估算的焦耳热不能作为严格能量分配恒等式，并指出撕裂/等离子体微观效应破坏该层流模型的条件。
